\section{Relation-aware compilation and bounded execution}
\label{ws:design}
\subsection{The program retains interactions, not edge arrays}
A graph maps each edge to a source and a destination. A message function reads
source, destination and edge fields; a reducer combines messages for each
destination. The result is
$$
y_i=\operatorname{finalize}\!\left(
\bigoplus_{e:d(e)=i}\operatorname{lift}
\big(f(x_{s(e)},z_i,a_e;\theta)\big)\right).
$$
The reducer declares an identity, state update and finalization. Duplicate
edges are separate contributions; empty rows follow the reducer's identity
and finalization convention. An optional node function consumes the reduced
value. The edge function determines what is computed, while the graph determines
which interactions exist.

For example, this program sums source features over a coordinate-dependent
radius relation. Positions and features are ordinary Torch tensors on the same
device; rows denote destinations and columns denote sources.
\begin{lstlisting}[language=Python]
import tiga as tg

class NeighborSum(tg.MessagePassing):
    reducer = tg.sum()
    def edge(self, src, dst, edge):
        return src.x

graph = tg.Graph.radius(positions, cutoff=0.1)
y = NeighborSum()(graph=graph, src={"x": x}, dst={})
\end{lstlisting}
The interface does not require an edge list, although a consumer may select an
explicit representation. Dense and triangular graphs similarly describe an
interaction domain. For causal attention, the destination supplies a query,
the source supplies a key/value, and an online-softmax reducer combines score
and value messages. Preserving the reducer state admits streaming evaluation
without a global score matrix, following the IO-aware principle exemplified by
FlashAttention \citep{flashattention}.

\subsection{Lowering preserves legality before choosing a schedule}
Tiga uses MLIR for its compiler representations \citep{mlir}. Domain IR retains
field roles, graph construction and message/reducer regions. Iter IR introduces
relation-coordinate traversal. Kernel IR assigns tiles and reduction work to
execution resources. Task/Storage IR describes dependencies, physical copies
and ownership where the selected path needs them. Tensor IR carries numerical
expressions alongside these layers; it is not a mandatory final stage.

Three properties constrain the transformations. First, field roles distinguish
a source ID from an edge position, even when both index floating-point arrays.
Second, reducer associativity permits regrouping and commutativity permits
reordering; arbitrary Python functions do not automatically satisfy either
property. Third, shape, dtype, layout and topology guards determine whether an
executable remains legal. Floating-point reduction changes require numerical
tolerances even when real-arithmetic semantics agree. NVIDIA lowering emits
serialized Triton IR for provider compilation \citep{triton}; native CPU paths
lower through LLVM. A provider interface alone does not establish a working
backend on another vendor's hardware.

\paragraph{Generated relations.}
For supported low-dimensional Euclidean radius workloads, a cell directory
limits traversal to neighboring cells. An exact distance predicate selects
accepted sources, whose fields can be consumed immediately by the reduction.
This avoids storing accepted-edge indices, distances and messages globally.
A custom metric needs a compatible spatial bound before that traversal is legal.

For exact kNN, a destination program scans candidate tiles and keeps a padded
top-k state, ordered by squared distance and source ID. If the smallest key
in a candidate tile cannot improve the retained state, sorting and merging
that tile are skipped; its distances are still computed. This is exact
selection pruning, not approximate neighbor search or a subquadratic algorithm.
Compatible new coordinates can rebind an existing executable without capture
and lowering. Coordinates remain live inputs: executable reuse never implies
reuse of previously selected neighbors.

\subsection{Paging separates capacity from kernel speed}
\label{ws:paging}
Explicit graphs larger than device memory use a native storage path. Topology
and source fields remain in a persistent store. The runtime loads a bounded
destination-row page, gathers its source fields on the host, stages the page
to CUDA and runs its reduction. It assembles the final output in a resident
device tensor. The capacity constraint is therefore
$$
M_{\mathrm{output}}+M_{\mathrm{active\ topology}}
+M_{\mathrm{staged\ fields}}+M_{\mathrm{temporary}}
\le B_{\mathrm{device}},
$$
where the budget applies to Tiga-managed device buffers, not other processes
or the full CUDA context. The output must still fit the device. Page topology
and field versions must agree; transfer must finish before consumption; and
storage cannot be recycled until its consumers finish.

The current implementation gathers features in page-edge order, potentially
repeating a source value many times. It also returns completed output pages
through host bytes before writing the resident result. These choices bound
temporary memory but can dominate runtime. The measured configuration is
serialized, without prefetch. Optional read-ahead does not implement a complete
disk--host--device overlap pipeline. Ordinary Torch tensors remain under Torch's
allocator; Tiga-managed paging uses native tensors rather than silently taking
ownership of arbitrary Torch storage.

\paragraph{Differentiation boundary.}
Supported Torch-facing programs interoperate with Torch autograd. CUDA EdgeNN
paths recompute local activations in a symbolic vector--Jacobian product (VJP);
some other compiled forwards replay Torch operations during backward. CPU
paging has a supported native VJP, but CUDA paged execution and the dense
streaming attention path evaluated here have no compiled backward. Derivatives
act on the selected relation, not discrete radius membership or kNN selection.
Appendix~\ref{ws:gradients} distinguishes the numerical checks from training
performance evidence.
